\section{Modelos de Linguagem (LLMs) na Odontologia}
\label{sec:ia-llms}

A convergência multifacetada entre IA generativa profunda (LLMs de trilhões de parâmetros), gêmeos digitais (DT) e ecossistemas da Internet das Coisas (IoT) estabelece um novo e incontornável paradigma interdisciplinar, desdobrando o espectro de atuação das máquinas muito além da mera análise de pixels e vóxeis \cite{mdpi2025digital}. 

Os Modelos de Linguagem de Grande Escala, como as iterações robustas e com arquitetura \textit{Mixture of Experts} (MoE), atenuam o ruído semântico existente nos intrincados relatórios médicos e anamneses orais. Eles conferem ao gêmeo digital odontológico, essencialmente numérico, uma "interface dialógica natural", possibilitando interações cognitivas bidirecionais entre o avatar do paciente e a equipe clínica.

\textbf{Geração procedimental e auditoria clínica:} Os LLMs traduzem vetores espaciais advindos de segmentações (como a contiguidade de uma lesão de furca no elemento 46) em laudos descritivos minuciosamente formatados para exigências periciais e convênios, isentando o cirurgião do custo burocrático, enquanto reduz as incongruências linguísticas no prontuário.

\textbf{Agentes de suporte à decisão (CDS):} Não se tratam de meros oráculos bibliográficos. Em frameworks de orquestração multiagente amparados por Retrieval-Augmented Generation (RAG), como os encontrados no ecossistema OpenCode, agentes especializados cruzam a ficha medicamentosa do paciente com toda a base do PubMed e do repositório ChEMBL simultaneamente. Eles emitem contraindicações farmacodinâmicas absolutas ou ajustes posológicos com base na idade metabólica aferida virtualmente \cite{technologysociety2025}.

\textbf{Tutoria contínua para educação em saúde:} Em hospitais e clínicas universitárias, gêmeos digitais com interface conversacional simulam comportamentos reativos complexos, viabilizando o ensino simulado de endodontia e o treinamento exaustivo sem qualquer desgaste ou risco a tecidos vivos, como ratificado pelos debates mais recentes na Nature BDJ \cite{naturebdj2026}.
